\sectioncentered*{Перечень условных обозначений}
\addcontentsline{toc}{section}{Перечень условных обозначений}
\label{sec:reduction}
\begin{flushleft}
	\setlength{\parindent}{1.25cm}
		\indent
	\indent API --- Application Programming Interface, программный интерфейс приложения;\\
	\indent БД --- база данных;\\
	\indent СУБД --- система управления базами данных;\\
	\indent ЛВС --- локальная вычислительная сеть;\\
	\indent ИПС --- информационно-поисковая система;\\
	\indent СИП --- система информационного поиска;\\
	\indent HTML --- HyperText Markup Language, язык гипертекстовой разметки;\\
	\indent HTTP --- HyperText Transfer Protocol, протокол передачи гипертекста;\\
	\indent JSON --- JavaScript Object Notation, текстовый формат обмена данными;\\
	\indent REST --- Representational State Transfer, архитектурный стиль взаимодействия компонентов;\\
	\indent SQL --- Structured Query Language, язык структурированных запросов;\\
	\indent TF --- Term Frequency, частота термина в документе;\\
	\indent IDF --- Inverse Document Frequency, инверсная документная частота;\\
	\indent TF-IDF --- Term Frequency --- Inverse Document Frequency, статистическая мера значимости термина;\\
	\indent VSM --- Vector Space Model, векторная модель;\\
	\indent NLP --- Natural Language Processing, обработка естественного языка;\\
	\indent NLTK --- Natural Language Toolkit, библиотека обработки естественного языка;\\
	\indent pgvector --- расширение PostgreSQL для хранения и поиска векторных представлений;\\
	\indent CSV --- Comma-Separated Values, формат с разделителями-запятыми;\\
	\indent DOCX --- формат документов Microsoft Word;\\
	\indent PDF --- Portable Document Format, переносимый формат документов;\\
	\indent RTF --- Rich Text Format, формат расширенного текста;\\
	\indent F-мера --- мера качества поиска, объединяющая точность и полноту;\\
	\indent P@k --- Precision at k, точность на первых k результатах;\\
	\indent watchdog --- библиотека Python для мониторинга файловой системы;\\
\end{flushleft}
